\MajorSection{إلى القارئ}

لقد اجترأ المترجم على أن يسمّي هذا التأليف، الذي يطمح إلى أن يسبق زمانه، «أنشودة الناموس الأعلى»؛ ولم يخف خطر اصطدام الاسم بصيغ غير مستساغة مثل «الثقافة العليا». والمبادئ التي تسوّغ هذه التسمية هي الآتية.

يقرر المؤلف أن السعادة والشقاء مقسومان وموزعان بالتساوي في العالم.

ويجعل تنمية النفس، مع المراعاة الواجبة للآخرين، الغاية الوحيدة والكافية للحياة الإنسانية.

ويرى أن العواطف والتعاطف و«عطية الشفقة الإلهية» هي أسمى متع الإنسان.

ويدعو إلى تعليق الحكم، مع الارتياب الملائم في «الوقائع، أكثر الخرافات عبثًا».

وأخيرًا، وإن بدا هدّامًا، فهو في جوهره يعيد البناء.

ولمزيد من التفاصيل عن القصيدة والشاعر، يُحال القارئ الفضولي إلى آخر المجلد.

ف ب

فيينا، تشرين الثاني ١٨٨٠.
